\chapter{The small-ratio weak-C2 opening is also admitted for all tests}
\label{component:p296e824cb97a}
\noindent\textbf{Component }\path{p296e824cb97a}\par
{\small\noindent Source: \path{R13/advances/weak_A_boundary/extension.tex}\par}
{\small\noindent SHA256: \path{296e824cb97a0923083aae8ef939c13988961e23d165a42f581ecd58ab6690b4}\par}
{\small\noindent Inventory use: WE1-WE11\par}
\medskip
\section{The small-ratio weak-C2 opening is also admitted for all tests}
\label{p296e824cb97a:sec:r10-weak-boundary}

\begin{theorem}
Retain the original physical C2 source, its exact stationary contact,
all strict source gaps, $n\ge0$, and the normalization $d=1$.
Suppose $B,z,y>0$ are infinitesimal,
$r,1-r,G/B,j,y/B,L/B$ are positive bounded units, and $H/y$ is
infinitesimal. If $a=z/B$ is positive infinitesimal, the entire
original parent is a finite tensor SOS at every independent
projective testing pair, for all real or complex states, including
signed charts and literal zero and infinite half-angles.
No comparison between source infinitesimal rates or between source
and testing rates is imposed. In particular $\delta_{\rm tot}>0$
is infinitesimal. This admission does not assert nonemptiness of
its intersection with the previous residual complement.
\end{theorem}

The original omitted event coefficients remain $\lambda=\beta=0$.
Here $B=H+L$, $h=H/B$, $t=y/B$, and the original contact is
\[
 (u_c,v_c;W_c,X_c,Y_c,Z_c)
 =\left(\sqrt z,\sqrt y;B/\sqrt{zy},-r\sqrt z,H/\sqrt y,1\right).
\]
No residual-complement premise is used. The proof extends the exact
six-contact full-parent Gram (WA3)--(WA8) and its finite separated
lift (WA10); it does not use the disproved global polynomial splitter.

\paragraph{Exact cancellation, with no $h/a$ restriction.}
Put $\Delta=\delta/B$, $M_*=M/B^2$, $\mu=m/a^2$ and
$\alpha_*=\alpha/B^2$. The unchanged source formulas give
\[
\begin{gathered}
 c_a=1+aBr,\quad S=1+B(1-h+ar),\\
 g=G/B=1+a(hB+r-1),\quad j=c_a-(1-h)/t,\\
 \begin{aligned}\Delta={}&a+(1+aB)+(1+aB)B+2aB(1+aB)r\\ &+a^3B^2r^2
       +2a^2Br+tB[1+a+aB+a^2Br(2-r)]\end{aligned},\\
 M_*=(1+aB)(1+ar)(h+t)-a(1+tB)(1-r)(1-h),\\
 n=\frac{(1+ar)c_aS-r\Delta}{2r\Delta},\qquad
 \mu=\frac{(1+tB)(1-r)c_aS}{2g\Delta},\\
 C=\frac{c_aM_*}{2(1-h)\Delta},\quad
 Ah=\frac{t c_a(1+ar)S}{2\Delta},\\
 D=\frac{(1+aB)(1+ar)(1-atBr)S}{2tj\Delta},\quad
 \alpha_* =\frac{(1+ar)M_*}{2(1-hB)\Delta}.       \end{gathered}\tag{WE1}
\]
All displayed denominators are units at the joint corner
$B=a=h=0$; $j>0$ a unit implies $t_0=\operatorname{st}(t)>1$.
The subscript $0$ denotes residue, and $0<r_0<1$.
Thus
\[
 n_0=\frac{1-r_0}{2r_0},\quad \mu_0=\frac{1-r_0}2,
 \quad C_0=\frac{t_0}2,\quad D_0=\frac1{2(t_0-1)},
 \quad (Ah)_0=(\alpha_*)_0=\frac{t_0}2.          \tag{WE2}
\]
These are positive; $m=a^2\mu$ is infinitesimal, $A$ is
unbounded, and $Ah^2,\sqrt A h,A^{-1/2}$ are infinitesimal.
All other retained coefficients are bounded. Consequently
$\delta_{\rm tot}=1/(1+n+m+C+A+D+\alpha)$ is positive
infinitesimal, as required in the original source normalization.

The only new potentially unbounded local coefficient is
$(b_2-\alpha_*)/a$, where
\[
 b_2=2a^2\mu h+C(1-h)^2+(Ah)h+(B^2\alpha_*+1/2)h^2.
\]
At $a=0$, with $d_*=1+B(1+t)$, the exact values are
\[
 C=\frac{h+t}{2(1-h)d_*},\quad
 Ah=\frac{t[1+B(1-h)]}{2d_*},\quad
 \alpha_* =\frac{h+t}{2(1-Bh)d_*}.
\]
Substitution gives $b_2=\alpha_*$ identically for arbitrary
$B,h,t$. Clearing the unit denominators in (WE1), the numerator
of $b_2-\alpha_*$ is therefore divisible by $a$ in the polynomial
ring. Its quotient has a unit denominator at $B=a=h=0$.
Hence
\[
 Q_*=(b_2-\alpha_*)/a\quad\hbox{is bounded}.       \tag{WE3}
\]
This is a joint regularity argument, so it permits arbitrary
$h/a$, including an unbounded ratio. The checker verifies the
exact vanishing, polynomial divisibility, and nonzero denominator
at the joint corner.

\paragraph{The complete local five-root Gram survives.}
Keep $\kappa=B/y$ and every coefficient and normal in
(WA3)--(WA8), but set $\ell=1/a$ and replace $R_2$ by
\[
 \widetilde R_2=aR_2=avW-uZ.
\]
The entire full-parent matrix is the exact congruence
\[
 \widetilde G=EGE,\qquad E=\operatorname{diag}(1,1,1,a^{-1},1),
 \qquad \widetilde{\mathcal R}=(\xi,q,R_1,\widetilde R_2,T).
                                                               \tag{WE4}
\]
Its newly rescaled coefficients are
\[
\begin{gathered}
 m/a^2=\mu,\qquad
 \beta_s/a^2=(Bd_4+\kappa c_4)/\kappa,\\
 J_s(U)/a=\frac{aBd_4-a_4\kappa^2+d_2+d_4U}{2\kappa},\qquad
 f_s/a=\tfrac12[Q_*-ac_4-aBd_4/\kappa],\\
 (-m+q_1U)/a=-a\mu+(nr+Bp_1)U,\\
 q_y/a=-k_z-\kappa p_y'-Bk_{uv}.                  \end{gathered}\tag{WE5}
\]
All are bounded, by (WE1)--(WE3). The other coefficients are
bounded as before; the complete second normal remains divided
by $\sqrt A$. Thus no first normal or hard-state term is omitted.
For independent infinitesimal or zero $u,v$ the complete residue is
\[
 \overline{\widetilde G}
 =\operatorname{diag}(n_0+1/2,1,C_0+D_0+1/2,\mu_0,t_0/2)
 \succ0.                                               \tag{WE6}
\]
In particular the previous $\xi,R_2$ cross entry becomes
$-a\mu$ and has zero residue. Applying the pure-source and
pure-testing finite lift (WA10) proves the full local statement.

\paragraph{Bounded tests with nonzero $u$ residue.}
Normalize $\widehat Y=\sqrt A Y$ without centering and use the
complete three-state residue (WR2), now with $m_0=0$ and (WE2).
For $U=u_0^2$, $V=v_0^2$, its exact determinant is
\[
 \det P_3=(1+U)(1+V)U[b_*U+c_*UV+d_*V^2],             \tag{WE7}
\]
where
\[
 b_*=\frac{n_0(2C_0+2D_0+1)}4,\quad
 c_*=\frac{C_0(2D_0+2n_0+1)}4,\quad
 d_*=\frac{C_0D_0(2n_0+1)}2.
\]
Every coefficient is positive. For $U>0$, the restriction to
$Z=0$ is positive definite from the positive $n_0,C_0$ packets.
The determinant therefore gives a positive definite complete
residue, with the additional hard $\widehat Y$ pivot
$(1+U)(1+V)$. The finite separated lift around $(u_0,v_0)$ applies.

\paragraph{The missing finite-$v$ and reciprocal-$v$ charts.}
For $u$ infinitesimal or zero and $v$ a bounded unit, use the five
roots
\[
 (S_*,R_*,X,\widehat Y,Z)=(aW,uW,X,\sqrt A Y,Z),\qquad
 Y=A^{-1/2}\widehat Y.
\]
Writing $D_u=1+u^2$, $D_v=1+v^2$, the following is the entire
original $F/4$, as a quadratic polynomial in those roots:
\[
\begin{split}
 nD_uD_vX^2&+\mu(vS_*-aX+auvY-auZ)^2
 +CD_v(R_*-Y)^2+D_uD_v\widehat Y^2\\
 &+D(R_*+uvX-Y-vZ)^2+\alpha D_u(vY-Z)^2\\
 &+\tfrac12\{(R_*+uvX-Y-vZ)^2+D_u(vY-Z)^2
               +D_v(X+uZ)^2-D_uD_vZ^2\}.             \end{split}\tag{WE8}
\]
The coefficients are polynomials in $u,v$. Their pure-source factors $a,\mu,n,C,D,\alpha,A^{-1/2}$ are bounded.
At $u_0=0$, $v_0\ne0$, the separate pivots on
$(S_*,X,\widehat Y)$ are
$\mu_0v_0^2,(n_0+1/2)(1+v_0^2),1+v_0^2$.
The remaining block on $(R_*,Z)$ is
\[
 \begin{pmatrix}
 C_0(1+v_0^2)+D_0+1/2&-(D_0+1/2)v_0\\
 -(D_0+1/2)v_0&D_0v_0^2
 \end{pmatrix},\qquad \det=C_0D_0v_0^4>0.           \tag{WE9}
\]
Here $C_0D_0-D_0/2-1/4=0$ exactly. Thus (WA10) applies
in $u,v-v_0$, including literal $u=0$; roots need not be
independent for a positive definite Gram to give the asserted SOS.

For $v$ unbounded or literally infinite, put $q_v=1/v$ and
use the polynomial homogeneous parent $q_v^2F(u,1/q_v)/4$.
In the same five roots it is exactly
\[
\begin{split}
 nD_u(1+q_v^2)X^2
 &+\mu(S_*-aq_vX+auY-auq_vZ)^2
 +C(1+q_v^2)(R_*-Y)^2+D_u(1+q_v^2)\widehat Y^2\\
 &+D(q_vR_*+uX-q_vY-Z)^2+\alpha D_u(Y-q_vZ)^2\\
 &+\tfrac12\{(q_vR_*+uX-q_vY-Z)^2+D_u(Y-q_vZ)^2\\
 &\hspace{20mm}+(1+q_v^2)(X+uZ)^2-D_u(1+q_v^2)Z^2\}.
                                                               \end{split}\tag{WE10}
\]
For independent infinitesimal or zero $u,q_v$ its full residue
is diagonal with pivots
\[
 (\mu_0,C_0,n_0+1/2,1,D_0)>0.                       \tag{WE11}
\]
Its coefficients are bounded for the same reason as (WE8).
The separated lift again applies, including $q_v=0$ literally.
The checker expands (WE8) and (WE10) back to all six original
event packets and the complete signed core, and checks both residues.

\paragraph{Remaining endpoints and exhaustive coverage.}
If $u$ is unbounded or infinite and $v$ bounded, (WR4) is still
positive definite with $m_0=0$: its positive coefficients are
$C_0,n_0,D_0+1/2,1/2$, and the hard pivot is $1+v_0^2$.
If $v$ is unbounded or infinite and $u$ is a bounded unit,
(WR5) is positive definite because $C_0u_0^2>0$,
$n_0+1/2>0$, and $D_0>0$, with hard pivot $1+u_0^2$.
If both are unbounded or infinite, use the full homogeneous
five-root identity (WR6). Its residue (WR7) is still strictly
positive: the $Q$ pivot becomes $1/2$, and the other nontrivial
block determinant remains $D_0(n_0+1/2)>0$.
All these identities have bounded coefficients also at $m_0=0$.

For two bounded tests, both zero residues use (WE4)--(WE6);
nonzero $u$ residue uses (WE7); zero $u$ and nonzero $v$ residue
uses (WE8)--(WE9). A single unbounded coordinate is covered by
the preceding paragraph, except for unbounded $v$ and zero $u$
residue, which is (WE10)--(WE11). Two unbounded coordinates use
(WR6)--(WR7). These cases exhaust every independent testing type.
Literal projective endpoints use the exact polynomial chart value.
Away from literal endpoints, restoration multiplies by a
pure-testing square $p^{-2}$, $q_v^{-2}$, or $(pq_v)^{-2}$.
No division is made at a literal endpoint. All four original
homogeneous state coordinates are retained; real and imaginary
parts give the complex-state result. The proof establishes the
full-parent all-tests admission on this source opening, without
a scalar-to-parent inference or a new physicality assumption.

\paragraph{The unit-ratio assumption can now be removed.}
Under the retained hypotheses, $H\in\Inf$ and $1-r$ is a
positive unit. Since $g=1-a(1-r-H)>0$, one has
$0<a<(1-r-H)^{-1}$, so $a$ is bounded. Every positive such
$a$ is either a unit, covered by (WA1), or infinitesimal, covered
here. Their union therefore gives the same all-tests theorem
without the additional $z/B$-unit assumption. This argument does
not cover a vanishing $1-r$ or any other omitted source gap.
